\clearpage
\section{유한 자동동형군과 아르틴의 고정체 정리}
\label{sec:artin-fixed-field}

\begin{learninggoal}
유한 자동동형군 $H$가 고정하는 체 $L^H$를 기준체로 삼으면 $L/L^H$가 갈루아 확장이 됨을 증명한다. 확장차수가 $|H|$와 같고 갈루아군이 정확히 $H$임을 보인다.
\end{learninggoal}

부분군에서 고정체로 가는 구성은 쉽다. 더 깊은 사실은 유한한 자동동형군을 고정체 위의 갈루아군으로 정확히 복원할 수 있다는 것이다.

\begin{lemma}[궤도다항식]
$H$가 체 $L$의 유한 자동동형군이고 $K=L^H$라 하자. $a\in L$의 $H$-궤도에서 서로 다른 원소들을
\[
  a_1,\dots,a_r
\]
라 하면
\[
  P_{a,H}(X)
  :=\prod_{i=1}^r(X-a_i)
\]
는 $K[X]$에 속하는 분리다항식이다. 따라서 $a$는 $K$ 위에서 분리적이고
\[
  [K(a):K]\le r\le |H|.
\]
\end{lemma}

\begin{proof}
임의의 $\tau\in H$는 궤도 $\{a_1,\dots,a_r\}$를 순열한다. 따라서 $\tau$를 계수에 적용해도 곱은 변하지 않으며, $P_{a,H}$의 모든 계수는 $H$의 원소들에 의해 고정된다. 즉 계수는 $K=L^H$에 속한다. 서로 다른 일차인수의 곱이므로 $P_{a,H}$는 분리다항식이고, $a$의 최소다항식은 이를 나누므로 결론이 따른다.
\end{proof}

\begin{theorem}[아르틴의 고정체 정리]
$L$이 체이고 $H\le\Aut(L)$가 유한부분군이라 하자. $K=L^H$라 두면
\[
  L/K
\]
는 유한 갈루아 확장이고
\[
  [L:K]=|H|,
  \qquad
  \Gal(L/K)=H.
\]
\end{theorem}

\begin{proofidea}
궤도다항식을 이용하면 모든 원소가 $K$ 위에서 분리적이고 차수가 $|H|$ 이하임을 알 수 있다. 유한 부분확장들의 차수가 균일하게 제한되므로 전체 확장도 유한하다. 마지막으로
\[
  H\subseteq\Aut_K(L),
  \qquad
  |\Aut_K(L)|\le[L:K]
\]
를 비교하면 모든 부등식이 등호가 된다.
\end{proofidea}

\begin{proof}
$n=|H|$라 하자. 앞의 보조정리에 의해 모든 $a\in L$는 $K$ 위에서 분리적 대수원소이고
\[
  [K(a):K]\le n.
\]
따라서 $L/K$는 분리 대수적 확장이다.

이제 $E/K$가 $L/K$ 안의 임의의 유한 부분확장이라고 하자. $E/K$는 유한 분리확장이므로 원시원 정리에 의해 $E=K(b)$인 $b\in E$가 존재한다. 그러므로
\[
  [E:K]=[K(b):K]\le n.
\]
즉 모든 유한 부분확장의 차수는 $n$ 이하이다.

가능한 차수들은 $1,2,\dots,n$ 가운데 있으므로 차수가 최대인 유한 부분확장 $E/K$를 택할 수 있다. 임의의 $a\in L$에 대해 $E(a)/K$도 유한 분리확장이고 그 차수는 $n$ 이하이다. $E\subseteq E(a)$이므로 최대성에 의해
\[
  [E(a):K]=[E:K].
\]
따라서 $E(a)=E$이고 $a\in E$이다. 모든 $a\in L$에 대해 성립하므로 $L=E$이며 $L/K$는 유한확장이다. 또한
\[
  [L:K]\le n.
\]

정의상 $H$의 모든 원소는 $K=L^H$를 고정하므로
\[
  H\subseteq\Aut_K(L).
\]
한편 유한확장에 대한 일반 부등식을 사용하면
\[
  n=|H|
  \le|\Aut_K(L)|
  \le[L:K]
  \le n.
\]
따라서 모든 부등식은 등호이고
\[
  [L:K]=n,
  \qquad
  \Aut_K(L)=H.
\]
자동동형 개수에 의한 갈루아 판정으로 $L/K$는 갈루아 확장이다.
\end{proof}

\begin{corollary}[전체 갈루아군의 고정체]
$L/K$가 유한 갈루아 확장이고 $G=\Gal(L/K)$라 하면
\[
  L^G=K.
\]
\end{corollary}

\begin{proof}
$F=L^G$라 두자. 아르틴의 정리에 의해
\[
  [L:F]=|G|.
\]
한편 $|G|=[L:K]$이고 $K\subseteq F\subseteq L$이므로 탑 법칙에서 $[F:K]=1$이다.
\end{proof}

\begin{corollary}
$L/K$가 유한 갈루아이고 $H\le\Gal(L/K)$이면
\[
  L/L^H
\]
는 유한 갈루아 확장이며
\[
  \Gal(L/L^H)=H,
  \qquad
  [L:L^H]=|H|.
\]
\end{corollary}

\begin{example}
$L=\Q(\sqrt2,\sqrt3)$와
\[
  H=\langle\sigma_2\rangle
  =\{\id,\sigma_2\}
\]
를 생각하자. $\sigma_2$는 $\sqrt2$의 부호를 바꾸고 $\sqrt3$을 고정하므로
\[
  L^H=\Q(\sqrt3).
\]
아르틴의 정리는
\[
  [L:\Q(\sqrt3)]=2=|H|
\]
와
\[
  \Gal(L/\Q(\sqrt3))=H
\]
를 동시에 준다.
\end{example}

\begin{remark}[무한확장에서 생기는 현상]
$H$가 무한이면 $L/L^H$가 대수적이라는 추가 가정 아래 여전히 갈루아 확장이 되지만, $H$가 전체 갈루아군과 같다고 보장되지는 않는다. 무한 갈루아 대응에서는 갈루아군에 크룰 위상을 주고 \emph{닫힌 부분군}만을 사용해야 한다. 이 주제는 다음 권에서 다룬다.
\end{remark}

\begin{checkpoint}
\begin{enumerate}[label=(\alph*)]
  \item 궤도다항식의 계수가 고정체에 들어가는 이유를 대칭다항식의 관점에서 설명하라.
  \item 아르틴의 정리에서 $L/K$가 유한임을 보이는 데 차수의 균일한 상계가 어떻게 사용되는지 설명하라.
  \item $H=\{\id\}$이면 아르틴의 정리가 어떤 자명한 사실을 말하는지 적어라.
\end{enumerate}
\end{checkpoint}
